% Phase 2 paper - IEEEtran conference template, 8-10 pages including references.
%
% THIS FILE IS A SKELETON. The %-comments under each heading are notes: the facts, numbers and citations
% each section needs. Write every paragraph yourself, in your own words, then delete the notes.
% Numbers: never type a result by hand. Use \R{key} (keys listed in numbers_index.md), so every number
% in the paper comes from results/ and is reproducible with the scripts.
%
% Build: bash paper/build_assets.sh  (copies tables, figures, numbers from results/), then upload the
% paper/ folder to Overleaf (or: cd paper && latexmk -pdf main.tex).

\documentclass[conference]{IEEEtran}
\usepackage[T1]{fontenc}
\usepackage{graphicx}
\usepackage{booktabs}
\usepackage{amsmath,amssymb}
\usepackage{url}
\usepackage[hidelinks]{hyperref}
\usepackage{xcolor}
\usepackage{cite}
\graphicspath{{figures/}}
\input{numbers}

\begin{document}

\title{TITLE: your own title}
% Notes: your Phase 1 title, or a sharper one now that results exist. Name the system and the claim.

\author{
\IEEEauthorblockN{Your Name (Roll No.)}
\IEEEauthorblockA{Project G1-XX \\ FAST-NUCES \\ email@nu.edu.pk}
% Brief, section 5.1: title page must include title, project code, and names, roll numbers and emails
% of all group members.
}

\maketitle

\begin{abstract}
% 150-250 words, one paragraph, in this order:
% - problem: heterogeneous robot teams + natural-language mission briefs; LLM planners (e.g. COHERENT)
%   are evaluated on success rates only, never against optimal allocation, and do not verify plans
% - approach: tool-using agent; formal model -> verified solver portfolio; cross-check by an independent
%   second reading; staged extraction + clarification for informal briefs
% - evaluation: 50-scenario benchmark (\R{bench/test} test scenarios x explicit/informal briefs),
%   model(s) used, base paper reproduced (COHERENT-style PEFA)
% - main numbers: agent feasible \R{<model>/agent:full/feasible/all}\% vs PEFA
%   \R{<model>/llm-pefa/feasible/all}\%; success@5; drop explicit->informal; ablation highlight
% - contribution sentence
\end{abstract}

\begin{IEEEkeywords}
multi-robot task allocation, large language models, LLM agents, heterogeneous robots, verification
\end{IEEEkeywords}

\section{Introduction}
% ~0.9 column. Points, in order:
% - motivation: mixed fleets (drones, rovers, mobile manipulators, fixed arms) at inspection /
%   maintenance sites; operators describe missions in words, not MILPs
% - LLM planners now allocate tasks from language: COHERENT \cite{liu2025coherent}, LaMMA-P
%   \cite{zhang2025lammap}, EMOS \cite{chen2025emos}, LiP-LLM \cite{obata2025lipllm}
% - the three gaps (from Phase 1), one sentence each with evidence:
%   G1 no comparison with optimal/classical allocation (COHERENT: success + steps only; LiP-LLM:
%      execution time not considered)
%   G2 no formal verification (COHERENT checks single actions; IROS'25 MILP builder 82% extraction
%      accuracy \cite{peng2025milp}; Nl2Hltl2Plan open-loop \cite{xu2025nl2hltl2plan})
%   G3 informal language barely tested (LiP-LLM 0.93 -> 0.44 on ambiguous instructions)
% - research questions RQ1-RQ3 (one per gap; reuse your Phase 1 wording, Table in B.11)
% - contributions (bulleted list, 3-4 items): benchmark + paired explicit/informal briefs with ground
%   truth; agent with switchable gap components; evaluation against exact/heuristic allocators and the
%   reproduced base paper on the same data; open code (repo link)

\section{Related Work}
% ~1 column, synthesis of the Phase 1 review, grouped by theme (not paper by paper):
% (a) LLM planners for multi-robot teams: COHERENT (PEFA loop, BEHAVIOR-1K \cite{li2022behavior1k}),
%     SMART-LLM \cite{kannan2024smartllm}, RoCo \cite{mandi2024roco}, centralised vs decentralised
%     \cite{chen2024scalable}, EMOS, VIKI-R \cite{vikir2025}, hierarchical planner \cite{kawabe2026hierarchical}
% (b) LLM + symbolic solver: LaMMA-P (PDDL), LiP-LLM (LP), Nl2Hltl2Plan (LTL), LLM+P \cite{liu2023llmp},
%     CaStL \cite{guo2025castl}
% (c) LLMs writing optimisation models: OptiMUS \cite{ahmaditeshnizi2024optimus}, Autoformulation
%     \cite{astorga2025autoformulation}, DRoC \cite{jiang2025droc}, AFL \cite{zhang2026afl}, MILP for MRTA
%     \cite{peng2025milp}
% (d) classical MRTA: taxonomies \cite{gerkey2004taxonomy,korsah2013taxonomy,nunes2017taxonomy};
%     Hungarian \cite{kuhn1955hungarian}; auctions \cite{lagoudakis2005auction,choi2009cbba}
% (e) agents: ReAct \cite{yao2023react}, Toolformer \cite{schick2023toolformer}, Reflexion
%     \cite{shinn2023reflexion}, self-consistency \cite{wang2023selfconsistency} (related to our cross-check),
%     SayCan \cite{ahn2022saycan}, Inner Monologue \cite{huang2022innermonologue}
% End with: what none of them does (the three gaps), and where this paper sits.

\section{Problem Formulation}
% ~0.8 column. Use docs/formal_model.md (rewrite in your notation):
% - instance: robots (mobility air/ground/fixed, capabilities, speed, payload, endurance, start),
%   tasks (location, required capabilities, payload, duration, optional deadline), precedence pairs,
%   standby rules, horizon H
% - plan = ordered route per robot; timing rule (start = max(arrival, predecessors finished))
% - constraints C1-C8 (coverage, capability, payload, reachability, standby, precedence/deadlock,
%   deadlines/horizon, endurance); objective makespan C_max
% - class: ST-SR-TA with cross-schedule dependencies \cite{korsah2013taxonomy}, NP-hard
% - input to the LLM methods: only the brief (explicit or informal); output: a plan; the checker
%   judges it against the hidden ground truth
% - hypotheses H1-H3 (one per RQ, from Phase 1 B.11)

\section{Proposed System}
% ~1.3 columns + Fig. 1 (full width).
% - Fig. 1: \includegraphics{architecture.pdf}; caption names the gap components and the baselines
% - agent loop: JSON {thought, action, args}, one tool per turn, max 12 steps, state (draft, model,
%   solution, checks, questions), trajectory logging
% - tools table (README "Architecture"): draft_model, submit_spec (+ edit language), cross-check,
%   solve, check + finish gate, ask_human, finish
% - G1 component (mrta/gaps/g1_portfolio.py): portfolio auction -> Hungarian -> MILP with time budget
%   min(60, max(2, 0.25 |T||R|)) s; returns best plan feasible on the agent's model
% - G2 component (g2_verify.py): sanity checks; independent second reading (different prompt),
%   fact-by-fact diff (tasks, robots, precedence, standby, horizon); agent must review before solving;
%   finish gate refuses a plan violating the model once
% - G3 component (g3_language.py): 3 passes (fleet/map; tasks; deadlines/precedence/standby) with a
%   reading guide; edits instead of retyping; ask_human <= 2 questions, simulated owner answers only
%   with brief lines naming the asked ids
% - robustness: parse-failure reminders, action aliases, fallbacks (portfolio on stored model; draft)
% - each component switchable (AgentConfig presets used in E5)

\section{Experimental Setup}
% ~0.8 column.
% - benchmark v2: Table (benchmark_summary.tex) - 50 scenarios, S20/M20/L10, dev 10 / test 40, fixed;
%   generation + witness + why (docs/benchmark.md); difficulty: random plans feasible
%   \R{bench/random-feasible/S}/\R{bench/random-feasible/M}/\R{bench/random-feasible/L}\%;
%   auction infeasible on \R{bench/auction-infeasible} scenarios
% - informal briefs: how written (docs/benchmark.md section 4), one short example pair
% - methods: classical (MILP 60 s, Hungarian, auction) on ground truth; LLM-direct; LLM-formulate;
%   COHERENT-style PEFA (state what was kept: assigner + executability feedback + 5-round history,
%   and what was not: separate executor LLMs, BEHAVIOR-1K); agent full; ablations
% - models: exact ids (e.g. qwen2.5:7b-instruct via Ollama, 16k context; gemini-2.5-flash),
%   temperature 0.2, seeds 0-2, max tokens; hardware (your laptop)
% - repeats: 3 for E1-E3, 1 for E5; all prompts in mrta/llm/prompts.py; no tuning on test

\section{Evaluation Methodology}
% ~0.5 column.
% - feasibility: plan passes all C1-C8 on the ground truth (same checker for every method)
% - gap to optimal: (makespan - ref)/ref, ref = best classical makespan (MILP optimum on
%   \R{classical/milp/feasible/S}, ... proved optimal counts from E1_classical)
% - success@5 = feasible and gap <= 5%
% - model fidelity: fact-level precision/recall/F1 and recall per fact type (mrta/spec_eval.py)
% - efficiency: LLM calls, tokens, latency, solver time, tool calls, cost
% - statistics: mean +- std over repeats; McNemar exact test \cite{mcnemar1947} for paired feasibility;
%   Wilcoxon signed-rank \cite{wilcoxon1945} for gaps; bootstrap 95% CI
% - base paper comparison: reproduced by us on the same scenarios/metrics (say so explicitly);
%   COHERENT's own BEHAVIOR-1K numbers are not comparable (different task and metric)

\section{Results}
% ~1.5 columns, organised BY GAP (brief, section 5.2):
% - Table: main_results_compact.tex (or main_results.tex if space)
% - Gap 1: agent vs PEFA / direct / formulate - \R{<model>/gap1/vs-llm-pefa/delta} pp, CI, p;
%   gap to optimal vs heuristics by tier (Fig. feasibility_by_tier); answer RQ1, H1 supported or not
% - Gap 2: violations by type (E4_gap2_violations_by_type.tex); what the cross-check caught
%   (E4_gap2_crosscheck.tex: F1 before -> after \R{<model>/gap2/f1_first} -> \R{<model>/gap2/f1_final});
%   answer RQ2
% - Gap 3: explicit vs informal (E4_gap3_explicit_vs_informal.tex, Fig. explicit_vs_informal);
%   recall by fact type (E4_gap3_recall_by_constraint.tex); answer RQ3
% - gap-resolution table (gap_resolution.tex) - REQUIRED in the paper
% - report negative or partial results honestly (brief: they are not penalised)

\section{Ablation and Comparative Analysis}
% ~0.7 column.
% - Table E5_ablation.tex + Fig. ablation: effect of removing G1, G2, G3 one at a time (p-values)
% - efficiency trade-off: E6_efficiency.tex + Fig. tokens_vs_feasibility (cost of the agent's extra
%   calls vs reliability gained); MILP time vs portfolio
% - second model (if run): does the ranking of methods hold?

\section{Failure Analysis}
% ~0.7 column. From results/failures/<model>/:
% - table of categories F1-F8 with counts (categories.md) for the agent and for PEFA
% - 2 success + 2 failure cases (cases/*.md): short trajectory excerpts (which step went wrong)
% - for each major category: likely cause + how it could be fixed (future work)

\section{Limitations and Threats to Validity}
% ~0.5 column. From docs/benchmark.md section 3:
% - generated benchmark, one domain, regular phrasing; all numbers present in briefs
% - auction feasible on most scenarios -> feasibility mostly measures reading the brief
% - capability guide written with knowledge of the phrase banks
% - simulated mission owner
% - base paper re-implemented, not its original code/simulator; no BEHAVIOR-1K
% - model size / one local model; temperature 0.2; 3 repeats
% - no physical robots; ROS 2 execution not evaluated

\section{Conclusion}
% ~0.3 column: one sentence per gap with the key number; future work (dynamic re-planning with
% mission events, ROS 2 / Gazebo execution, larger and human-written brief sets, other domains).

\bibliographystyle{IEEEtran}
\bibliography{refs}

\end{document}
